% !TEX root = ../nc-cubic.tex
\subsection{周期付き生成対象による抽象残余圏の認識}
\label{cog:sec:recognition}
方程式を復元できても、その圏への比較は別に必要である。以下では比較核を作るために
一対象の内部シフト間の代数を調べる。係数体は $\C$ とし、$N=n+2,\ n\ge2$ とする。
$J$ を dg 自己同値、$P$ を対象とし、合成を内部シフトで移送して
\[
 B_j=\operatorname{RHom}(P,J^jP),\qquad B=\bigoplus_{j\in\mathbb Z}B_j
\]
と置く。第一次数をコホモロジー、第二次数を内部次数とする。
ここでの周期の仮定は $J^3\simeq[2]$ の比較と両立する
$D=\C[T,T^{-1}]$ 線型の二重次数付き dg 模型をこの軌道代数が持つこととする。
$|T|=(2,-3)$、微分は $(1,0)$ で、$T$ は中心的である。
具体的には各 $B_j$ の模型が復元する射複体とその合成（高次整合を含む）が
元の軌道のものに擬同値で、$T$ 倍が指定の周期を与えることを要求する。
単に三乗周期が三角圏上の自己関手の同型として存在するだけでは、この仮定を課した
ことにはならない。以下で生成対象は compact 対象の厚生成の意味で用いる。

\begin{theorem}[三重 Massey 積と可換三次の実比較]\label{cog:thm:recognition}
冪等完備な smooth proper $\C$ 線型 dg 圏 $\mathcal C$ について、次は同値である。
\begin{enumerate}
\item $\mathcal C$ は滑らかな三次 $n$ 次元超曲面の $\Ku$ に dg Morita 同値である。
\item ある $J,P$ と上記の周期模型が存在し、$P,JP,J^2P$ は厚生成する。
$U=H^{1,-1}(B)$ は次元 $N$ で、積により
\[
               \bigwedge U\otimes D\xrightarrow{\ \sim\ }H^{*,*}(B)
\]
となる。また下記の自己三重 Massey 積の三次式 $F_P$ が $\PP(U)$ で
滑らかな超曲面を定める。
\end{enumerate}
$N=3m+1$ なら $J=S^{-1}[m+1]$、$N=3m+2$ なら $J=S[-m]$ を第二条件で
指定できる。$3\mid N$ の場合は $J$ の存在を条件に含める。
比較先は $X=(F_P=0)$ であり、比較は生成対象間の双加群で構成される。
\end{theorem}

$a\in U$ を表すサイクル $\theta$ を取り、$dh=\theta^2$ と解く。
$\theta^2$ は外積の仮定で完全であり、$h$ の二重次数は $(1,-2)$ である。
符号の規約を固定して
\[
                 F_P(a)=-T^{-1}[\theta h+h\theta]\in\C
\]
と定める。$H^{1,-2}=0$ なので三重 Massey 積の不定性は零である。
従ってサイクルと原始元の選択によらない。重みを保つ線型な代表選択と
微分への線型切断を取れば $\theta$ は $a$ に線型、$h$ は二次であり、
これは三次式である。周期の非零スカラー倍は方程式の零点を変えない。

\begin{proof}
まずこの外積型の軌道代数を三次式から比較できることを示す。
$A=\bigwedge U\otimes D$ とする。微分を収縮し $B$ の厳密単位的極小
$A_\infty$ 模型を $A$ に移す
\cite[Sections~3.3, 3.5 and~7]{COGKellerInfinity}。
内部次数の三つの剰余類の代表で体上のサイクルと境界に補空間を選び、
他の次数には $T$ 倍で移すので、収縮は $D$ 線型で二重次数を保てる。
単位は代表に含める。$m_d$ の次数は $(2-d,0)$ である。

$A$ の $D$ 相対、正規化した Hochschild 複体で、arity $d$、
値の次数差 $(2-d,0)$ のコホモロジーを計算する。
対角の Koszul 双加群分解の長さ $d$ の生成空間は
$\operatorname{Sym}^d U$ で、各生成元の内部二重次数は $(d,-d)$
（分解の長さを別に数える）である。実際、符号付き包絡代数では左右の
奇生成元の差 $\delta_i$ も外積変数となり、対角はこれらで割った商である。
一変数では $\delta_i$ 倍を反復する分解の像と核が等しく、これを全変数で
tensor すれば主張の自由分解を得る。対角を値とする Hom に移すと $\delta_i$
の作用が零なので微分は消える。bar 分解への比較は奇生成元の懸垂が偶で
あるため、長さ $d$ の線型入力を対称化して送る写像である。
具体的に bar 微分の隣り合う二入力の積は順序を逆にした項と相殺され、
端の二項が上の $\delta_i$ 微分を与える。従ってコホモロジー類は
線型入力上の対称化で読み取れる（体の標数は零）。

この次数で可能な値 $T^b\bigwedge^a U$ は
\[
                     a+2b=2,\qquad a+3b=d
\]
に従う。$d=3$ では $a=0,b=1$ であり、障害空間は
$T\operatorname{Sym}^3 U^*$ である。$d\ge4$ では $a=6-2d<0$ で全て消える。
$m_3$ のクラスの対角入力での値は、transfer の二つの木の式で
上の自己三重 Massey 積になる（全体符号は両模型で同一）。
三次式が一致する二模型では従って $m_3$ の差は境界である。
一般に $m_2,\ldots,m_{d-1}$ を一致させた後、Stasheff 恒等式から
二つの $m_d$ の差は Hochschild cocycle である。その類が零なら
arity $d-1$ の正規化した cochain による、線型部が恒等の座標変更で消せる。
上の消滅によりこれを全 arity まで反復できる。各 tensor 長では有限回で安定するので、
bar 余代数上の可逆な比較を得る。これは $m_d$ 自体が $d\ge4$ で全て零で
あるとの仮定を必要としない。

次に可換三次 $F(x)=c(x,x,x)$ の模型を計算する。$c$ は対称三重形式、
$x_i$ は $U$ の基底の双対座標である。剰余体の次数付き Koszul 安定化の
微分と閉演算子を
\[
 d_K=\sum_i x_i\iota_i+T\sum_i c(e_i,x,x)\epsilon_i,\qquad
 \theta_a=\iota_a-T\sum_i c(a,e_i,x)\epsilon_i
\]
と取る。$\epsilon_i$ は外積、$\iota_i$ は収縮であり、
$d_K^2=TF$、$[d_K,\epsilon_i]=x_i$ となる。
$\epsilon^I\theta^J$ の順序基底での自己射の微分は $x_i$ の Koszul 微分であるため、
係数を原点に落とした外積代数が積を含むコホモロジーである。
これは \cite[Theorems~4.1, 5.9]{ROGDyckerhoff}の安定化模型の斉次三次の場合である。
ここでは三次項の係数と符号を直接検算できる。実際に
\[
 \theta_a^2=-T c(a,a,x),\qquad
 h_a=-T\sum_i c(a,a,e_i)\epsilon_i,\qquad
 [d_K,h_a]=\theta_a^2,\qquad \theta_a h_a+h_a\theta_a=-TF(a).
\]
これで先の規約は $F_P=F$ を与える。

滑らかな射影三次では錐の臨界軌跡は原点であり、この安定化の全内部シフトが
次数付き特異圏を厚生成する。次数付きでの生成を明記する。各偏微分の乗法は
行列因子化上、微分行列をその変数で偏微分したホモトピーにより零となる。
Jacobian イデアルは有限余長なので十分高い各 $x_i$ の冪も零作用である。
任意の次数付き特異圏の対象にこれらの冪の Koszul 複体を tensor すれば、
各乗法の cone が分裂するので元の対象は得られた対象の直和因子になる。
全 cone は有界導来圏で有限長の次数付きコホモロジーを持つから、
剰余体のシフトで構成される。周期により必要な内部シフトは三代表に尽きる。
次数付き Orlov 比較 \cite[Theorems~2.5, 3.10]{NFMOrlovSingularities} と
標準の Serre 公式 $S=J^{-N}[N]$ から必要性および $J$ の表示が従う。

逆に、第二条件の $B$ と $F_P$ の Koszul 安定化の軌道代数で $U,T$ を同定する。
上の帰納的比較によって二重次数を保つ $A_\infty$ 擬同型の列が得られる。
内部次数 $j-i$ の複体を対象 $i$ から $j$ への射とし、$0\le i,j\le2$ に
制限すれば、実際の三対象の full dg 部分圏の間の比較になる。
bar 型のグラフ双加群で表し、源側の Yoneda と標的生成対象への導来 tensor を
合成する。両側で三対象が厚生成するのでこれと逆比較は全圏の dg Morita 同値を
与える。よって対象の存在と有限次数の積・Massey 積の照合が十分である。
\end{proof}

\begin{corollary}[Serre 環で復元した標的に対する生成対象判定]
\label{cog:cor:serre-witness}
五次元で $J=S^{-1}[3]$ を固定する。命題\ref{nft:prop:jacobian-necessary}の
二次関係の積分条件を満たし、さらに全自然変換環が一次生成でこれらの二次式を
正則列とする完全交叉であると仮定する。
積分した滑らかな三次式を $F$ とする。この抽象圏が $\Ku(V(F))$ で実現されるには、
前定理の周期・外積型の厚生成対象が存在し、その $F_P$ が線型座標変換と
非零スカラー倍で $F$ に一致することが必要十分である。
\end{corollary}
\begin{proof}
十分性は定理の双加群と方程式の座標同型のグラフの合成である。
必要性では可換標的の Koszul 生成対象と周期を実同値で引き戻せる。
複数の積分の選択は命題\ref{gcd:prop:integrability-symmetrizer}により
滑らかな場合は射影同型を与える。
\end{proof}
指定二因子 Sklyanin の滑らかな可換比較点と偶符号の非零係数点では、
既構成の双加群の逆でこの生成対象を得る。偶符号の滑らかな零係数点では
三重積の階数 $34$ により滑らかな三次を与える生成対象は存在しない。
定理は候補対象を知らない場合にも適用できる存在判定であるが、
Serre 環のベクトル空間と乗法だけから対象や周期模型を取り出す構成は含まない。
